
%----------------- LISTA DE ABREVIATURAS E SIGLAS--------------------------------------
\siglalista{ABNT}{Associação Brasileira de Normas Técnicas}

%Para usar uma dada sigla ABC ao longo do texto, use \acrfull{ABC} se quiser apresentar a sigla e sua definição.
%Se quiser apresentar apenas a sigla, use \acrshort{ABC}.






%-----------------SÍMBOLOS---------------------------------------------------------------

\simbololista{G}{\ensuremath{G}}{Resultante do peso próprio (ação permanente)}
\simbololista{gamma\_med}{\ensuremath{\gamma_{med}}}{Valor médio do peso específico do material}
\simbololista{V}{\ensuremath{V}}{Volume do elemento estrutural}
\simbololista{F\_d}{\ensuremath{F_d}}{Valor de cálculo da ação combinada}
\simbololista{gamma\_gi}{\ensuremath{\gamma_{gi}}}{Coeficientes de majoração das ações permanentes}
\simbololista{F\_Gi,k}{\ensuremath{F_{Gi,k}}}{Valores característicos das ações permanentes}
\simbololista{gamma\_q1}{\ensuremath{\gamma_{q1}}}{Coeficiente de majoração da ação variável principal}
\simbololista{gamma\_qj}{\ensuremath{\gamma_{qj}}}{Coeficientes de majoração das ações variáveis secundárias}
\simbololista{F\_Q1,k}{\ensuremath{F_{Q1,k}}}{Valor característico da ação variável principal}
\simbololista{psi\_0j}{\ensuremath{\psi_{0j}}}{Fatores de combinação das ações variáveis secundárias}
\simbololista{F\_Qj,k}{\ensuremath{F_{Qj,k}}}{Valores característicos das ações variáveis secundárias}
\simbololista{m}{\ensuremath{m}}{Número de ações permanentes}
\simbololista{n}{\ensuremath{n}}{Número de ações variáveis secundárias}
\simbololista{gamma\_qw}{\ensuremath{\gamma_{qw}}}{Coeficiente de majoração da ação do vento}
\simbololista{gamma\_q2}{\ensuremath{\gamma_{q2}}}{Coeficiente de majoração da ação variável secundária}
\simbololista{psi\_0,q2}{\ensuremath{\psi_{0,q2}}}{Fator de combinação da ação variável secundária}
\simbololista{X\_d}{\ensuremath{X_d}}{Valor de cálculo da resistência da madeira}
\simbololista{X\_k}{\ensuremath{X_k}}{Valor característico da resistência da madeira}
\simbololista{k\_mod}{\ensuremath{k_{mod}}}{Coeficiente de modificação da madeira}
\simbololista{gamma\_w}{\ensuremath{\gamma_w}}{Coeficiente de ponderação da resistência}
\simbololista{k\_mod,1}{\ensuremath{k_{mod,1}}}{Fator de modificação associado à classe de carregamento}
\simbololista{k\_mod,2}{\ensuremath{k_{mod,2}}}{Fator de modificação associado à classe de umidade}
\simbololista{E\_0,ef}{\ensuremath{E_{0,ef}}}{Valor efetivo do módulo de elasticidade}
\simbololista{E\_0,med}{\ensuremath{E_{0,med}}}{Valor médio do módulo de elasticidade}
\simbololista{sigma\_c0,d}{\ensuremath{\sigma_{c0,d}}}{Tensão normal de cálculo na compressão paralela às fibras}
\simbololista{M\_sd}{\ensuremath{M_{sd}}}{Momento fletor solicitante de cálculo}
\simbololista{W\_c}{\ensuremath{W_c}}{Módulo de resistência da seção transversal comprimida}
\simbololista{f\_b,d}{\ensuremath{f_{b,d}}}{Resistência de cálculo à flexão}
\simbololista{sigma\_r0,d}{\ensuremath{\sigma_{r0,d}}}{Tensão normal de cálculo na tração paralela às fibras}
\simbololista{W\_t}{\ensuremath{W_t}}{Módulo de resistência da seção transversal tracionada}
\simbololista{I}{\ensuremath{I}}{Momento de inércia da seção transversal}
\simbololista{y\_c}{\ensuremath{y_c}}{Distância da linha neutra ao bordo comprimido}
\simbololista{y\_t}{\ensuremath{y_t}}{Distância da linha neutra ao bordo tracionado}
\simbololista{b}{\ensuremath{b}}{Largura da seção transversal retangular}
\simbololista{h}{\ensuremath{h}}{Altura da seção transversal retangular}
\simbololista{R}{\ensuremath{R}}{Raio da seção transversal circular}
\simbololista{D}{\ensuremath{D}}{Diâmetro da seção transversal circular}
\simbololista{tau\_d}{\ensuremath{\tau_d}}{Valor de cálculo da tensão de cisalhamento}
\simbololista{V\_d}{\ensuremath{V_d}}{Esforço cortante de cálculo}
\simbololista{S}{\ensuremath{S}}{Momento estático da parte da seção transversal}
\simbololista{f\_v0,d}{\ensuremath{f_{v0,d}}}{Resistência de cálculo ao cisalhamento paralelo às fibras}
\simbololista{A}{\ensuremath{A}}{Área da seção transversal}
\simbololista{tau\_max}{\ensuremath{\tau_{max}}}{Tensão máxima de cisalhamento}
\simbololista{delta\_inst}{\ensuremath{\delta_{inst}}}{Flecha instantânea}
\simbololista{delta\_inst,G\_i,k}{\ensuremath{\delta_{inst,G_i,k}}}{Flechas instantâneas decorrentes das ações permanentes}
\simbololista{delta\_inst,Q\_1,k}{\ensuremath{\delta_{inst,Q_1,k}}}{Flecha devido à ação variável principal}
\simbololista{delta\_inst,Q\_j,k}{\ensuremath{\delta_{inst,Q_j,k}}}{Flechas decorrentes das ações variáveis secundárias}
\simbololista{psi\_1,j}{\ensuremath{\psi_{1,j}}}{Fator de redução para condição frequente}
\simbololista{delta\_fin}{\ensuremath{\delta_{fin}}}{Flecha final}
\simbololista{delta\_fin,G\_i,k}{\ensuremath{\delta_{fin,G_i,k}}}{Flechas finais associadas às ações permanentes}
\simbololista{delta\_fin,Q\_j,k}{\ensuremath{\delta_{fin,Q_j,k}}}{Flechas finais associadas às ações variáveis}
\simbololista{delta\_creep}{\ensuremath{\delta_{creep}}}{Flecha associada aos efeitos da fluência}
\simbololista{phi}{\ensuremath{\varphi}}{Coeficiente de fluência da madeira}
\simbololista{sigma\_d,f}{\ensuremath{\sigma_{d,f}}}{Valor máximo de cálculo da tensão atuante de flexão}
\simbololista{f\_d,f}{\ensuremath{f_{d,f}}}{Valor de cálculo da resistência à flexão}
\simbololista{M}{\ensuremath{M}}{Momento fletor de cálculo}
\simbololista{W}{\ensuremath{W}}{Módulo de resistência da seção transversal}
\simbololista{sigma\_d,f,x}{\ensuremath{\sigma_{d,f,x}}}{Tensão máxima de cálculo devido à flexão na direção do eixo x}
\simbololista{sigma\_d,f,y}{\ensuremath{\sigma_{d,f,y}}}{Tensão máxima de cálculo devido à flexão na direção do eixo y}
\simbololista{sigma\_d,n}{\ensuremath{\sigma_{d,n}}}{Valor de cálculo da tensão normal de tração}
\simbololista{f\_d,t}{\ensuremath{f_{d,t}}}{Resistência de cálculo à tração paralela às fibras}
\simbololista{f\_d,c}{\ensuremath{f_{d,c}}}{Resistência de cálculo à compressão paralela às fibras}
\simbololista{Q\_g,k}{\ensuremath{Q_{g,k}}}{Valor característico da força cortante máxima devido às ações permanentes}
\simbololista{g}{\ensuremath{g}}{Carga permanente distribuída por unidade de comprimento}
\simbololista{L}{\ensuremath{L}}{Vão do elemento estrutural}
\simbololista{Q\_q,k}{\ensuremath{Q_{q,k}}}{Valor característico da força cortante máxima reduzida}
\simbololista{P}{\ensuremath{P}}{Carga concentrada aplicada}
\simbololista{q}{\ensuremath{q}}{Carga distribuída}
\simbololista{a}{\ensuremath{a}}{Distância da carga concentrada ao apoio}
\simbololista{e}{\ensuremath{e}}{Excentricidade ou comprimento efetivo}
\simbololista{delta\_g,k}{\ensuremath{\delta_{g,k}}}{Valor característico da flecha máxima devida às ações permanentes}
\simbololista{E\_M,ef}{\ensuremath{E_{M,ef}}}{Módulo de elasticidade efetivo do material}
\simbololista{delta\_q,k}{\ensuremath{\delta_{q,k}}}{Valor característico da flecha máxima devida à ação variável}
\simbololista{b}{\ensuremath{b}}{Distância da carga concentrada ao apoio mais próximo}
\simbololista{M\_r,k}{\ensuremath{M_{r,k}}}{Valor característico do momento fletor máximo}
\simbololista{q\_r}{\ensuremath{q_r}}{Carga distribuída equivalente}
\simbololista{a\_r}{\ensuremath{a_r}}{Comprimento de distribuição da carga concentrada}
\simbololista{L\_r}{\ensuremath{L_r}}{Vão do tabuleiro}
\simbololista{k\_def}{\ensuremath{k_{def}}}{Fator para avaliação da deformação por fluência}
\simbololista{E\_mean,fin}{\ensuremath{E_{mean,fin}}}{Valor médio final do módulo de elasticidade}
\simbololista{G\_mean,fin}{\ensuremath{G_{mean,fin}}}{Valor médio final do módulo de cisalhamento}
\simbololista{K\_ser,fin}{\ensuremath{K_{ser,fin}}}{Valor final do módulo de deslizamento}
\simbololista{E\_mean}{\ensuremath{E_{mean}}}{Valor médio do módulo de elasticidade}
\simbololista{G\_mean}{\ensuremath{G_{mean}}}{Valor médio do módulo de cisalhamento}
\simbololista{K\_ser}{\ensuremath{K_{ser}}}{Módulo de deslizamento}
\simbololista{psi\_2}{\ensuremath{\psi_2}}{Fator para o valor quase-permanente da ação}
\simbololista{k\_def,1}{\ensuremath{k_{def,1}}}{Fator de deformação para o elemento 1 de madeira}
\simbololista{k\_def,2}{\ensuremath{k_{def,2}}}{Fator de deformação para o elemento 2 de madeira}
\simbololista{D\_min}{\ensuremath{D_{min}}}{Diâmetro mínimo da longarina}
\simbololista{M\_d}{\ensuremath{M_d}}{Momento fletor de cálculo}
\simbololista{Q\_d}{\ensuremath{Q_d}}{Esforço cortante de cálculo}
\simbololista{p\_f}{\ensuremath{p_f}}{Probabilidade de falha da estrutura}
\simbololista{X}{\ensuremath{\mathbf{X}}}{Vetor n-dimensional de variáveis aleatórias do sistema}
\simbololista{f\_X}{\ensuremath{f_X}}{Função densidade de probabilidade conjunta}
\simbololista{G}{\ensuremath{G}}{Equação do estado limite}
\simbololista{hat\{p\}\_f}{\ensuremath{\hat{p}_f}}{Estimativa da probabilidade de falha}
\simbololista{n\_s}{\ensuremath{n_s}}{Número de amostras}
\simbololista{I\_f}{\ensuremath{I_f}}{Função indicadora de falha (1 se falha, 0 caso contrário)}
\simbololista{g}{\ensuremath{g}}{Função estado limite (margem de segurança)}
\simbololista{beta}{\ensuremath{\beta}}{Índice de confiabilidade}
\simbololista{Phi^{-1}}{\ensuremath{\Phi^{-1}}}{Função de distribuição acumulada inversa da distribuição normal padrão}
\simbololista{P}{\ensuremath{P}}{Matriz de permutações para amostragem por hipercubo latino}
\simbololista{R}{\ensuremath{R}}{Matriz de números aleatórios uniformes}
\simbololista{S}{\ensuremath{S}}{Matriz de amostras normalizadas}
\simbololista{s\_ij}{\ensuremath{s_{ij}}}{Elemento da matriz S}
\simbololista{F^{-1}\_Xj}{\ensuremath{F_{X_j}^{-1}}}{Inversa da função de distribuição acumulada da variável Xj}
\simbololista{x\_ij}{\ensuremath{x_{ij}}}{Amostra da variável aleatória Xj}
\simbololista{mu\_g}{\ensuremath{\mu_g}}{Média da margem de segurança}
\simbololista{sigma\_g}{\ensuremath{\sigma_g}}{Desvio padrão da margem de segurança}
\simbololista{mu\_S}{\ensuremath{\mu_S}}{Média da solicitação}
\simbololista{sigma\_S}{\ensuremath{\sigma_S}}{Desvio padrão da solicitação}
\simbololista{mu\_R}{\ensuremath{\mu_R}}{Média da resistência}
\simbololista{sigma\_R}{\ensuremath{\sigma_R}}{Desvio padrão da resistência}
%Para usar um dado símbolo SIMB ao longo do texto, use \gls{SIMB}.
